problem:  
Let \((M^n, g, f)\) be a complete, noncompact, **steady gradient Ricci soliton** satisfying \(\mathrm{Ric}(g) = \nabla^2 f\) with positive Ricci curvature and bounded curvature.  
It is known that in all explicit examples (e.g. the Bryant soliton), the scalar curvature \(R(x)\) decays asymptotically like \(c\, r(x)^{-1}\), where \(r(x)\) denotes the distance from a fixed point.  

**Question:**  
Does there exist any complete, nonflat, positively curved, steady gradient Ricci soliton \((M, g, f)\) for which the scalar curvature decays strictly faster than that—i.e.  
\[
R(x) = o(r(x)^{-1}) \quad \text{as } r(x)\to\infty?
\]  
Equivalently, is the \(O(r^{-1})\) decay rate of the scalar curvature a universal lower bound for all positive steady Ricci solitons, or can genuinely faster-decaying steady solitons exist?  

If such solitons exist, what can be said about their asymptotic geometric structure (cylindrical, conical, or otherwise)? If they cannot exist, can one prove a rigidity statement asserting that \(R(x)\sim c\,r(x)^{-1}\) asymptotically for all such solitons?

Why is it a "good" problem:  
This problem asks whether the known scalar curvature decay rate of the Bryant soliton exhausts all possible behaviors among positively curved steady Ricci solitons. It directly addresses a **structural rigidity question**: is there a universal decay threshold that constrains all steady Ricci solitons? This connects curvature decay, asymptotic geometry, and uniqueness phenomena within Ricci soliton theory.  

The question sits at the crossroads of **geometric analysis** and **asymptotic geometry**: proving nonexistence of faster-decaying solitons would establish a quantitative rigidity principle analogous to canonical asymptotic limits in other geometric PDEs, while an existence result would open an unexpected new family of steady soliton asymptotics. Its statement is simple but its resolution would require profound understanding of the interplay between elliptic estimates, maximum principles, and large-scale geometric structure, making it a deep and conceptually central research problem.